% ============================================================
% Methodology section draft.
% \todo{...} marks values to fill from your study; do not submit with them.
% Section numbers and the code module names in comments map to
% docs/METHODOLOGY_GUIDE.md and docs/RUNBOOK.md.
% ============================================================

\section{Threat Model}
\label{sec:threat}

We consider an attacker who interacts with a deployed LLM only through its
public chat interface or API, with default system prompts and no access to
model weights. The attacker's goal is to obtain content that the model's
usage policy prohibits. Unlike optimization-based attacks
\citep{zou2023universal, chao2023jailbreaking}, our attacker uses no special
tooling: they simply write their request in Banglish, as millions of users do
daily. This makes the attack cheap, transferable across models, and hard to
distinguish from legitimate use. We additionally study a guarded deployment,
in which input and output are screened by a guard model
\citep{inan2023llamaguard, zeng2024shieldgemma}.

\section{The BanglishJail Benchmark}
\label{sec:dataset}

\paragraph{Harm taxonomy.}
We adopt the semantic categories of HarmBench \citep{mazeika2024harmbench}
for comparability and add five categories specific to the Bangladeshi
context: communal and religious incitement, election misinformation,
mobile-financial-service fraud, fake job and visa scams targeting migrant
workers, and gender-based harms (Table~\ref{tab:taxonomy}). We exclude
chemical, biological, radiological and nuclear weapons, content sexualizing
minors, and requests targeting real private individuals.

\begin{table}[t]
\centering
\small
\begin{tabular}{llrr}
\toprule
Group & Category & Harmful & Benign \\
\midrule
Standard & Cybercrime and hacking & \todo{n} & \todo{n} \\
 & Illegal activities & \todo{n} & \todo{n} \\
 & Harassment and hate & \todo{n} & \todo{n} \\
 & Misinformation & \todo{n} & \todo{n} \\
 & Self-harm & \todo{n} & \todo{n} \\
 & Fraud and scams & \todo{n} & \todo{n} \\
\midrule
Bangladesh- & Communal / religious incitement & \todo{n} & \todo{n} \\
specific & Election misinformation & \todo{n} & \todo{n} \\
 & Mobile financial service fraud & \todo{n} & \todo{n} \\
 & Job and visa scams & \todo{n} & \todo{n} \\
 & Gender-based harms & \todo{n} & \todo{n} \\
\midrule
 & Total seeds & \todo{N} & \todo{M} \\
\bottomrule
\end{tabular}
\caption{Harm taxonomy and number of seeds per category. Each seed has five
parallel versions.}
\label{tab:taxonomy}
\end{table}

\paragraph{Parallel versions.}
Each seed is written in five parallel versions with identical intent:
English (\textsc{en}), Bangla script (\textsc{bn}), standard Banglish
(\textsc{std}), noisy Banglish (\textsc{noisy}) and Banglish--English
code-mixed text (\textsc{cm}). Standard Banglish follows a single phonetic
convention close to the Avro keyboard mapping. Noisy Banglish applies
spelling patterns that we derived from \todo{N} public Banglish social-media
comments: vowel dropping (\textit{kemon}~$\rightarrow$~\textit{kmn}), sound
merging (\textit{bhalo}~$\rightarrow$~\textit{valo}) and conventional short
forms (\textit{ki}~$\rightarrow$~\textit{k}). We quantify noise as the
character edit distance between the noisy and standard versions, normalized
by length (mean \todo{x}). Code-mixed versions replace 20--40\% of content
words with English; on a manually tagged sample of \todo{n} prompts the
Code-Mixing Index \citep{gamback2014measuring} averages \todo{x}.

\paragraph{Authoring and quality control.}
\todo{N} native Bangla speakers authored the seeds following a written
guideline (Appendix~\todo{A}). Most seeds (\todo{x}\%) are native-authored;
the remainder are adapted from MultiJail \citep{deng2024multilingual} and
HarmBench \citep{mazeika2024harmbench} to enable comparison with prior work.
Each version was checked by a second annotator for meaning preservation and
naturalness (1--5); versions failing either check were revised or removed.
Agreement on a doubly-checked 20\% sample was $\kappa=\todo{x}$
\citep{cohen1960coefficient}.

\paragraph{Benign controls and splits.}
To measure over-refusal, we include \todo{n} benign controls that resemble
harmful requests but are harmless, typically protective (e.g., how to
recognize a fraudulent mobile-banking SMS). The final benchmark contains
\todo{N} harmful seeds and \todo{M} benign controls ($\times 5$ versions $=$
\todo{T} prompts). We split \emph{seeds} (not prompts) into a test set
(70\%) and a defense-training set (30\%), stratified by category, so that no
seed appears in both. A datasheet \citep{gebru2021datasheets} is provided in
the appendix.

\section{Evaluation Setup}
\label{sec:setup}

\paragraph{Models.}
We evaluate \todo{K} models: \todo{list proprietary models with exact
version identifiers}, open-weight models \todo{list}, served with vLLM
\citep{kwon2023efficient}, and \todo{Bangla-focused model(s)}. All queries
were made between \todo{dates}. We use each model's default system prompt,
greedy decoding where the API allows it, and a maximum output length of
\todo{512} tokens (\todo{4096} for models with built-in reasoning, whose
reasoning tokens count toward the limit). We report the decoding settings
of every model in Appendix~\todo{B}.

\paragraph{Attack conditions.}
Our main condition (\textsc{direct}) sends each version of each seed as a
single-turn request. We further run four conditions on subsets of the test
set: (i) a \emph{noise sweep}, perturbing 0, 25, 50, 75 and 100\% of the
words of standard Banglish to obtain a dose--response curve; (ii) a
\emph{response-language} condition requesting the answer in English or in
Banglish; (iii) \emph{template} attacks wrapping each version in \todo{n}
published jailbreak templates; and (iv) automated attacks with PAIR
\citep{chao2023jailbreaking} and TAP \citep{mehrotra2024tree}, using the
HarmBench implementations with an attacker instructed to write in Banglish
and a budget of \todo{20} queries.

\section{Scoring}
\label{sec:scoring}

\paragraph{Two-axis labels.}
A model that does not refuse a Banglish request may simply have
misunderstood it. We therefore label each response for harmful seeds as
\textsc{refusal} (R), \textsc{misunderstood} (M), \textsc{partial} (P) or
\textsc{harmful} (H), and responses to benign controls as
\textsc{over-refusal} (R), \textsc{misunderstood} (M) or \textsc{answered} (A).
Requests declined by the provider's API-level safety system are labelled R.

\paragraph{Metrics.}
The \emph{attack success rate} (ASR) is the fraction of harmful-seed
responses labelled H (strict) or H or P (lenient). \emph{Comprehension} is
the fraction not labelled M. Because only on-topic harmful responses count as
successful attacks, a high ASR cannot be explained by misunderstanding. For
benign controls we report the \emph{over-refusal rate}.

\paragraph{LLM judge and human validation.}
We label all responses with an LLM judge (\todo{judge model}) using a
rubric that mirrors the human guideline. The judge always receives the
English reference of the seed, so that it knows the intended request; for
non-English responses it also receives an English translation produced by
the same model. Items the judge declines to grade are labelled by humans.
Two annotators independently labelled a stratified sample of \todo{n}
responses (balanced over models and versions), and an adjudicator resolved
disagreements. Inter-annotator agreement was $\kappa=\todo{x}$, and
judge--human agreement was $\kappa=\todo{x}$ overall, ranging from
\todo{x} (\todo{version}) to \todo{y} (\todo{version}) across versions
(Table~\todo{x}). Human labels replace judge labels wherever available.

\paragraph{Statistical analysis.}
We report 95\% confidence intervals from a bootstrap over seeds (2{,}000
resamples). As every seed appears in all versions, we compare each version
with English using McNemar's test \citep{mcnemar1947note} on paired outcomes
per model, with Holm correction for multiple comparisons
\citep{holm1979simple}. To estimate the overall effect of version while
accounting for seed difficulty, we fit a mixed-effects logistic regression,
$\text{harmful} \sim \text{version} \times \text{model} + (1 \mid \text{seed})$,
using lme4 \citep{bates2015fitting}.

\section{Guard-Model Evaluation}
\label{sec:guards}

We screen every prompt (input moderation) and every harmful or answered
response (output moderation) with \todo{list guard models}. We report recall
on harmful seeds and the false-positive rate on benign controls, per version.

\section{Mechanistic Analysis}
\label{sec:mech}

For the open-weight models we analyze why safety fails to transfer.
First, we compare the number of tokens per prompt across versions.
Second, following \citet{arditi2024refusal}, we compute a refusal direction
as the difference between the mean last-token residual-stream activations of
\todo{n} harmful and \todo{n} harmless English instructions, selecting the
layer whose projection best separates held-out instructions (largest
Cohen's $d$, among the first 80\% of layers). We then project the
activations of every version of every seed onto this direction, and compute
the layer-wise cosine similarity between each version and its English
counterpart. Third, as a causal test, we add the refusal direction to the
residual stream (scale $\alpha=\todo{x}$) while the model processes harmful
Banglish, and measure whether refusals are restored.

\section{Defenses}
\label{sec:defenses}

We evaluate three defenses on the held-out test seeds.
\textbf{D1 (normalization before moderation):} Banglish input is converted
to Bangla script with IndicXlit \citep{madhani2023aksharantar}, keeping
English words unchanged, or translated to English by an LLM, before the
guard model screens it.
\textbf{D2 (safety fine-tuning):} we fine-tune open-weight models with LoRA
\citep{hu2022lora} on the defense-training split, pairing harmful prompts in
all five versions with refusals in the matching language, and benign
controls with helpful reference answers to avoid over-refusal, mixed with
general instruction data in a 1:1 ratio.
\textbf{D3 (activation steering):} the intervention from
Section~\ref{sec:mech}, applied at inference time.
For each defense we report ASR, over-refusal, helpfulness on benign
Banglish requests, performance on \todo{a general benchmark}, added latency,
and robustness to an adaptive attacker who knows the defense
\citep{nasr2025attacker}, e.g., by choosing spellings that defeat the
transliterator.

\section{Ethical Considerations}
\label{sec:ethics}

The study was approved by \todo{ethics board, approval number}. Annotators
were adults who gave informed consent, were paid \todo{rate}, worked in
sessions of at most 90 minutes, and could skip items or stop at any time.
We excluded the most severe harm categories. Harmful prompts and model
outputs are not released publicly; the benchmark is available under gated
access with a usage agreement. We disclosed our findings to \todo{providers}
on \todo{date}, \todo{n} days before submission.
